\documentclass[10pt,a4paper]{amsart}
\usepackage[margin=19mm]{geometry}
\usepackage{amsmath,amssymb,mathtools}
\usepackage[scr=rsfs,cal=euler]{mathalfa}
\usepackage{xfrac}
\usepackage[unicode,hidelinks,bookmarks=false]{hyperref}
\newcommand{\ie}{i.e.\ }
\newcommand{\cf}{cf.\ }
\newcommand{\resp}{respectively\ }
\newcommand{\Subsection}{\subsection}
\providecommand{\nobreakdash}{\nobreak\hbox{-}}
\setlength{\emergencystretch}{2em}
\multlinegap0pt
\usepackage{amssymb}
\usepackage{mathtools}
\newtheorem{proposition}{Proposition}
\DeclareMathOperator{\Ann}{Ann}\DeclareMathOperator{\ann}{Ann}
\DeclareMathOperator{\Ext}{Ext}
\DeclareMathOperator{\Hom}{Hom}
\DeclareMathOperator{\Ht}{ht}
\DeclareMathOperator{\grade}{grade}
\title{Annihilator of $Ext^i_R(R/I,R)$}
\author{Mohsen Asgharzadeh}
\date{Source: arXiv:2601.20596v3 (CC BY 4.0)}
\begin{document}
\maketitle

\noindent\emph{Source and licence.} Annihilator of $Ext^i_R(R/I,R)$. Version arXiv:2601.20596v3 (CC BY 4.0), distributed by arXiv under the Creative Commons Attribution 4.0 International licence, http://creativecommons.org/licenses/by/4.0/, as recorded in the arXiv licence metadata for this exact version and shown on the abstract page https://arxiv.org/abs/2601.20596v3. Full licence notice: \texttt{licenses/arxiv-2601.20596-CC-BY-4.0.txt}.\par\smallskip

\noindent\textit{Editorial scope.} Counted unit: the article's proposition cl (source0.tex lines 1471--1480). The article's complete original proof of it is delivered with it (source0.tex lines 1482--1523, one \texttt{proof} environment of 42 lines). No mathematics is written, repaired or re-derived: the statement and the proof are the article's own lines, in the article's own notation, and no retained line is substituted or re-flowed.  No further source-local context is needed: every cross-reference of every retained line resolves inside the two delivered spans.  Nothing was omitted from any retained span, so the package carries no omission record.  No retained line contains a citation, so the wrapper carries no bibliography and no bibliography entry.  No retained line contains a figure environment, an image-inclusion command or drawing code, so no image is delivered and the package declares no asset dependency.  Editorial formatting only, in addition: an AMS \texttt{amsart} preamble that restates the article's own declarations and macros used by the retained lines, namely the environments \texttt{proposition} on separate counters, together with the standard packages those lines need.  The retained environments print in this document as Proposition 1 in order of appearance, whereas the article prints the same items with the numbering of its own full document.  The complete native source of the article as distributed in the arXiv e-print for this exact version is bundled unchanged as \texttt{sources/193536/source0.tex}.
\begin{proposition}\label{cl}
	Let $I$ be a trace ideal containing a regular element. Then
	\(
	\Ann\!\big(\Ext^1_R(R/I, R)\big) \supseteq \mathfrak{C}_R.
	\)
	In particular, if $\mathfrak{C}_R=\mathfrak{m}$ and $\Ht(I)=1$, then
	\(
	D(I)=\mathfrak{m}.
	\)
\end{proposition}

\begin{proof}
	Consider the short exact sequence
	\(
	0 \longrightarrow I \longrightarrow R \longrightarrow R/I \longrightarrow 0.
	\)
	Since $I$ contains a regular element, every homomorphism $R/I \to I$ is zero; hence
	\(
	\Hom_R(R/I,I)=0.
	\)
	Applying $\Hom_R(-,R)$ gives an exact sequence
	\[
	0 \longrightarrow R \longrightarrow I^\ast \longrightarrow \Ext^1_R(R/I,R) \longrightarrow 0,
	\]
	where $I^\ast=\Hom_R(I,R)$.
Because $I$ is a trace ideal, we have $I^\ast=\Hom_R(I,I)$. Moreover, for any ideal $I$ containing a regular element,
	\(
	\Hom_R(I,I)\subseteq \overline{R}
	\)
	inside the total ring of fractions. Therefore
	\[
	\Ext^1_R(R/I,R)\cong I^\ast/R
	= \frac{\Hom_R(I,I)}{R}
	\subseteq \frac{\overline{R}}{R}.
	\]
	By definition of the conductor,
	\(
	\mathfrak{C}_R = \Ann_R(\overline{R}/R),
	\)
	so $\mathfrak{C}_R$ annihilates every submodule of $\overline{R}/R$. Hence
	\(
	\Ann_R\!\big(\Ext^1_R(R/I,R)\big) \supseteq \mathfrak{C}_R.
	\)
For the final statement, assume $\mathfrak{C}_R=\mathfrak{m}$ and $\Ht(I)=1$. Since $I$ contains a regular element, $\grade(I,R)=1$, and therefore
	$\Ext^1_R(R/I,R)\neq 0$. We have
	\[
	\mathfrak{m}=\mathfrak{C}_R \subseteq \Ann_R\!\big(\Ext^1_R(R/I,R)\big) \subseteq R.
	\]
	Because $\Ext^1_R(R/I,R)\neq 0$ and $R$ is local, its annihilator must be a proper ideal; hence it equals $\mathfrak{m}$. By definition,
	\(
	D(I)=\Ann_R\!\big(\Ext^1_R(R/I,R)\big)=\mathfrak{m}.
	\)
\end{proof}

\end{document}
